\documentclass[11pt,a4paper]{article}
\usepackage[utf8]{inputenc}
\usepackage[T1]{fontenc}
\usepackage[french]{babel}
\usepackage{amsmath,amssymb}
\usepackage[dvipsnames]{xcolor}
\usepackage{booktabs}
\usepackage[margin=2cm]{geometry}
\usepackage{hyperref}

\title{Synthèse des schémas testés\\\large bas Mach et hypersonique --- flux de chaleur pariétal}
\author{Vincent Delmas}
\date{\today}

\begin{document}
\maketitle

\section*{Méthode}

Quatre critères d'acceptation : (1) préservation du champ de vitesse sur le tourbillon
de Gresho ; (2) convergence de la fluctuation de densité en fonction du Mach, au moins
$\mathrm{Ma}^1$ ; (3) robustesse sur Sedov triangulaire \emph{et} hexaédrique ;
(4) flux de chaleur pariétal sur le demi-cylindre à Mach~27, quad \emph{et} tri.

Tous les résultats ci-dessous proviennent de \texttt{test\_low\_mach\_compar\_resume/},
une campagne unique lancée avec un binaire unique, pour que les chiffres soient
comparables entre eux.

\section*{1. Un piège de lecture : le critère 2 ne se lit jamais seul}

C'est le résultat méthodologique le plus important de cette campagne.

\begin{table}[h]\centering
\begin{tabular}{lcc}
\toprule
schéma & pente Mach (quad) & vortex Gresho (quad) \\
\midrule
\texttt{two\_wave}            & $3{,}53$ & $4{,}12\cdot10^{-1}$ \\
\texttt{multi\_point}         & $1{,}70$ & $4{,}12\cdot10^{-1}$ \\
\texttt{modified\_three\_wave}& $1{,}70$ & $4{,}12\cdot10^{-1}$ \\
\texttt{three\_wave}          & $1{,}87$ & $2{,}72\cdot10^{-1}$ \\
\texttt{WIP}                  & $1{,}41$ & $1{,}54\cdot10^{-2}$ \\
\texttt{WIP2\_NOLM}           & $1{,}12$ & $\mathbf{1{,}26\cdot10^{-3}}$ \\
\bottomrule
\end{tabular}
\end{table}

Les quatre schémas qui « gagnent » le critère~2 ont en réalité \textbf{détruit le
tourbillon} : $40\,\%$ d'erreur sur la vitesse, contre $0{,}13\,\%$ pour
\texttt{WIP2\_NOLM}. Pour \texttt{two\_wave}, $L_2(\rho)$ tombe à
$1{,}28\cdot10^{-15}$ --- le zéro machine d'un écoulement au repos, pas une
convergence. Un schéma assez dissipatif rend ce critère arbitrairement bon.

\subsection*{1.1 La version Voronoï : une anticorrélation parfaite}

L'étude a été refaite sur un maillage de Voronoï polygonal (2500 cellules, même
nombre que le quad $50\times50$), avec cette fois l'\textbf{énergie cinétique
intégrée} comme mesure de la préservation du tourbillon, et non la RMSE d'un
profil relevé le long d'une ligne :

\begin{table}[h]\centering
\begin{tabular}{lccc}
\toprule
schéma & $E_{kin}$ retenue (Voronoï) & $E_{kin}$ (quad) & pente Mach (Voronoï) \\
\midrule
\texttt{multi\_point\_pressure}& $\mathbf{98{,}69\,\%}$ & $99{,}94\,\%$ & $0{,}96$ \\
\texttt{WIP2\_NOLM}            & $\mathbf{98{,}69\,\%}$ & $99{,}94\,\%$ & $0{,}97$ \\
\texttt{WIP}                   & $89{,}20\,\%$ & $99{,}60\,\%$ & $1{,}15$ \\
\texttt{three\_wave}           & $37{,}10\,\%$ & $51{,}71\,\%$ & $1{,}75$ \\
\texttt{modified\_three\_wave} & \textcolor{BrickRed}{$0{,}000\,\%$} & $0{,}000\,\%$ & $1{,}71$ \\
\texttt{multi\_point}          & \textcolor{BrickRed}{$0{,}000\,\%$} & $0{,}000\,\%$ & $1{,}75$ \\
\texttt{two\_wave}             & \textcolor{BrickRed}{$0{,}000\,\%$} & $0{,}000\,\%$ & $\mathbf{3{,}26}$ \\
\bottomrule
\end{tabular}
\end{table}

Les deux classements sont \textbf{exactement inverses} l'un de l'autre : rangés
par pente croissante on obtient \texttt{multi\_point\_pressure},
\texttt{WIP2\_NOLM}, \texttt{WIP}, puis le groupe
\texttt{modified\_three\_wave} / \texttt{three\_wave} / \texttt{multi\_point},
et enfin \texttt{two\_wave} --- c'est-à-dire l'ordre décroissant de l'énergie
cinétique retenue. \texttt{two\_wave} réalise la meilleure pente de toute la
campagne ($3{,}26$) en ayant conservé exactement \emph{zéro} énergie. Le
classement sur Voronoï est par ailleurs identique à celui sur quad : le maillage
polygonal ne réordonne rien, il confirme.

\emph{Convention} : toutes les pentes de ce document sont obtenues par moindres
carrés sur les quatre nombres de Mach (\texttt{analyze.py}), et non entre les
deux derniers points seulement --- les deux définitions diffèrent sensiblement
dès que la courbe n'est pas une droite parfaite.

\paragraph{Un avertissement sur la métrique elle-même.} La colonne « vortex
Gresho » du tableau précédent est une RMSE relevée par sonde le long d'une
ligne. Sur un maillage irrégulier, cette sonde \textbf{surestime l'erreur d'un
facteur $\sim20$} : elle affichait une dégradation $\times27$ de
\texttt{WIP2\_NOLM} entre quad et Voronoï ($1{,}26\cdot10^{-3}$ vers
$3{,}42\cdot10^{-2}$) qui n'existe pas --- l'énergie cinétique intégrée montre
$99{,}94\,\%$ vers $98{,}69\,\%$, soit une perte qui passe de $0{,}065\,\%$ à
$1{,}31\,\%$, réelle mais modeste. Pour comparer des topologies de maillage
différentes, il faut une grandeur intégrale, pas une sonde.

Le même mécanisme a été mesuré indépendamment côté lagrangien : \texttt{classic},
de la même famille à vitesse nodale que \texttt{multi\_point}, dissipe $99{,}996\,\%$
de l'énergie cinétique du vortex de Gresho en $t=0{,}2$. La cohérence entre les deux
solveurs confirme qu'il s'agit bien d'un défaut de schéma et non d'un artefact.

\section*{2. Les sept schémas de référence}

\begin{table}[h]\centering
\begin{tabular}{lcccc}
\toprule
& vortex (q/t) & Sedov tri/hex & flux chaleur quad & flux chaleur tri \\
\midrule
\texttt{multi\_point}          & 4{,}12e-1 / 3{,}99e-1 & ok / ok & 6{,}99e-4 & 1{,}67e-3 \\
\texttt{three\_wave}           & 2{,}72e-1 / 1{,}30e-2 & ok / ok & 4{,}91e-4 & 1{,}52e-3 \\
\texttt{two\_wave}             & 4{,}12e-1 / 3{,}99e-1 & ok / --- & 3{,}23e-2 & 5{,}30e-2 \\
\texttt{modified\_three\_wave} & 4{,}12e-1 / 3{,}99e-1 & ok / --- & 2{,}93e-3 & 6{,}49e-3 \\
\texttt{multi\_point\_pressure}& 1{,}27e-3 / 1{,}79e-3 & \textcolor{BrickRed}{crash} & \textbf{4{,}45e-4} & \textcolor{BrickRed}{crash} \\
\texttt{WIP}                   & 1{,}54e-2 / 1{,}22e-2 & ok / ok & 6{,}99e-4 & 3{,}50e-4 \\
\texttt{WIP2\_NOLM}            & \textbf{1{,}26e-3 / 1{,}80e-3} & ok / ok & 6{,}98e-4 & \textbf{3{,}46e-4} \\
\bottomrule
\end{tabular}
\end{table}

\texttt{WIP2\_NOLM} est le seul à passer les quatre critères. \texttt{multi\_point\_pressure}
(Del Grosso, Barsukow, Loubère \& Maire 2026) est le meilleur sur le flux quad et
excellent en bas Mach --- il est \emph{asymptotic-preserving} par construction --- mais
il s'arrête sur \texttt{sedov\_tri}, \texttt{sedov\_hex} et le demi-cylindre tri.
Comme \texttt{sedov\_hex} est hexaédrique, ce n'est pas une affaire de type de maille
mais de choc fort, ce qui correspond au positionnement que les auteurs donnent
eux-mêmes à leur solveur : le multi-point à \emph{vitesse} nodale est celui qu'ils
décrivent comme \emph{carbuncle-free} en hypersonique.

\section*{3. Le terme de diffusion additionnel $\epsilon$}

\subsection*{3.1 Cinq formulations}

Les quatre candidats de \texttt{tex/wip.tex} plus deux tentatives sans capteur.
Constat préalable : \textbf{le critère~2 ne discrimine pas ce terme} --- les quatre
capteurs donnent $6{,}00$ à $6{,}03\cdot10^{-8}$ à $\mathrm{Ma}=10^{-4}$. C'est donc
un degré de liberté disponible pour le flux de chaleur.

\begin{table}[h]\centering
\begin{tabular}{llcc}
\toprule
variante & $\epsilon$ & flux quad & flux tri \\
\midrule
\texttt{WIP2\_NOLM} & Ducros nodal, $4a_p$ & 6{,}98e-4 & \textbf{3{,}46e-4} \\
\texttt{WIP2\_EJUMP}& saut de pression nodal & 4{,}63e-4 & 1{,}56e-3 \\
\texttt{WIP2\_EDIV} & $-\nabla\!\cdot\!v$ nodal & \textbf{4{,}23e-4} & 2{,}43e-3 \\
\texttt{WIP2\_EMAX} & $\max$, porte $\mathrm{Ma}^2$ & 4{,}11e-3 & 5{,}55e-3 \\
\texttt{WIP2\_EPOS} & $\max(0,-\Delta v_n,\sqrt{|\Delta p|/\bar\rho})$, par face & 1{,}07e-3 & 4{,}08e-3 \\
\texttt{WIP2\_EPOS2}& $\max(0,-\Delta v_n,|\Delta p|/(\bar\rho\bar a))$, par face & 2{,}36e-3 & 3{,}14e-3 \\
\bottomrule
\end{tabular}
\end{table}

Deux enseignements. D'abord, \textbf{le capteur de Ducros déjà en place est le meilleur
des six sur tri}, d'un facteur $4{,}5$ sur son suivant --- ce n'était pas un choix par
défaut. Ensuite, \textbf{les deux formulations par face sont battues sur les deux
maillages} par les quatre capteurs nodaux, ce qui confirme la robustesse supérieure des
capteurs nodaux.

À noter que \texttt{EPOS2} égale pourtant la référence en bas Mach ($5{,}94\cdot10^{-8}$
contre $6{,}03\cdot10^{-8}$) sans aucun paramètre : l'échec ne se voit que sur le flux
pariétal. Sa première version \texttt{EPOS} plafonnait à $1{,}06\cdot10^{-7}$, parce que
$\sqrt{|\Delta p|/\bar\rho}$ a la bonne dimension mais la mauvaise échelle en Mach
($\Delta p = O(h)$ indépendamment du Mach dans ce régime) ; la conversion par impédance
$|\Delta p|/(\bar\rho\bar a) = O(h\,\mathrm{Ma})$ corrige cela.

\subsection*{3.2 L'amplitude gouverne le compromis quad/tri}

Le capteur de Ducros \emph{inchangé}, à trois amplitudes :

\begin{table}[h]\centering
\begin{tabular}{lccc}
\toprule
amplitude & flux quad & flux tri & Sedov tri (RMSE $\rho$) \\
\midrule
$4a_p$ (\texttt{WIP2\_NOLM}) & 6{,}98e-4 & \textbf{3{,}46e-4} & 6{,}37e-1 \\
$2a_p$ (\texttt{WIP2\_ED2})  & 5{,}48e-4 & 7{,}23e-4 & 5{,}63e-1 \\
$1a_p$ (\texttt{WIP2\_ED1})  & \textbf{4{,}63e-4} & 7{,}35e-4 & \textbf{4{,}99e-1} \\
\bottomrule
\end{tabular}
\end{table}

Les tendances sont monotones et \textbf{opposées entre quad et tri}. Baisser
l'amplitude améliore simultanément le flux quad et la précision Sedov, et dégrade le
flux tri. Il n'existe donc pas de réglage satisfaisant les deux maillages : c'est un
compromis structurel, pas un défaut de formulation. À amplitude égale ($1a_p$), la
\emph{forme} compte aussi : Ducros donne $7{,}35\cdot10^{-4}$ sur tri contre
$2{,}43\cdot10^{-3}$ pour $-\nabla\!\cdot\!v$.

\section*{3bis. Le maillage hybride quad + Voronoï, ou comment discriminer}

Le critère 4 posait un problème de lisibilité : \textbf{sur le maillage quad,
presque tous les schémas se superposent à la référence LAURA} et rien ne les
sépare. Un maillage hybride --- couche limite structurée en quadrangles sur
$r\in[1;1{,}04]$, lointain en Voronoï, conformes à l'interface, 9103 mailles ---
les sépare nettement. Il transpose directement le maillage \texttt{tri}, qui est
lui-même un hybride quad + triangles de 8326 mailles.

\begin{table}[h]\centering
\begin{tabular}{lccc}
\toprule
RMSE du flux de chaleur vs LAURA & quad & tri & \textbf{hybride} \\
\midrule
\texttt{WIP2\_NOLM}            & $6{,}98\cdot10^{-4}$ & $3{,}46\cdot10^{-4}$ & $\mathbf{3{,}90\cdot10^{-4}}$ \\
\texttt{WIP}                   & $6{,}99\cdot10^{-4}$ & $3{,}50\cdot10^{-4}$ & $4{,}11\cdot10^{-4}$ \\
\texttt{three\_wave}           & $4{,}91\cdot10^{-4}$ & $1{,}52\cdot10^{-3}$ & $7{,}06\cdot10^{-4}$ \\
\texttt{multi\_point}          & $6{,}99\cdot10^{-4}$ & $1{,}67\cdot10^{-3}$ & $1{,}23\cdot10^{-3}$ \\
\texttt{multi\_point\_pressure}& $4{,}45\cdot10^{-4}$ & \textcolor{BrickRed}{plante} & $2{,}06\cdot10^{-3}$ \\
\texttt{modified\_three\_wave} & $2{,}93\cdot10^{-3}$ & $6{,}49\cdot10^{-3}$ & $5{,}79\cdot10^{-3}$ \\
\texttt{two\_wave}             & $3{,}23\cdot10^{-2}$ & $5{,}30\cdot10^{-2}$ & $5{,}10\cdot10^{-2}$ \\
\bottomrule
\end{tabular}
\end{table}

\noindent
\textbf{\texttt{WIP2\_NOLM} est le meilleur schéma du maillage hybride}, et il y
est meilleur que sur quad. Le classement hybride suit celui du \texttt{tri} et
non celui du quad, ce qui est cohérent puisque l'hybride en est la transposition.

Deux faits méritent d'être relevés. \texttt{multi\_point\_pressure}, impeccable
sur quad, développe sur l'hybride un \textbf{pic au point d'arrêt à $0{,}0155$,
près du double de LAURA} --- un défaut de type carbuncle, exactement le point
faible attendu d'une pression nodale face à un choc fort, et c'est lui qui
explique son RMSE. Et il \emph{tourne jusqu'au bout} sur l'hybride alors qu'il
\emph{plante} sur \texttt{tri} : les polygones sont plus cléments que les
triangles, comme du côté lagrangien.

\paragraph{Une précaution d'interprétation.} L'interface est placée à
$r_{bl}=1{,}04$, reprise du maillage \texttt{tri}. Le choc de culot (décollement
de l'ordre de $0{,}15\,R$ à Mach~27, donc $r\approx1{,}15$) tombe donc dans la
région de \emph{Voronoï}, exactement comme il tombe dans la région triangulaire
du maillage \texttt{tri}. Le bloc structuré ne sert qu'à la couche limite et au
flux pariétal. Une dégradation éventuelle du flux ne pourrait donc pas être
attribuée à la paroi sans vérifier d'abord le traitement du choc par les
polygones, ce que trancherait un calcul à $r_{bl}=1{,}2$.

\section*{4. Changer d'advection pour se passer de $\epsilon$}

\subsection*{4.1 Les variantes ZB existantes}

\texttt{ZB\_ARMD\_LPP}, \texttt{ARMDU}, \texttt{ARMDMAT}, \texttt{ARMDUMAT} :
toutes crashent sur \texttt{sedov\_tri} et donnent $5{,}1\cdot10^{-3}$ à
$3{,}3\cdot10^{-1}$ sur le flux de chaleur, soit un à trois ordres de grandeur
au-dessus de la famille WIP. Elles hybrident le flux multi-D nodal et le flux 1D de
face avec $w = 0{,}5$ codé en dur, alors qu'\texttt{AMISO} --- la seule advection
robuste de cette famille --- utilise $w = (\min_f A_f/A_f)/3$, au plus un tiers et
pondéré par la géométrie.

\subsection*{4.2 Advection multi-D ré-hybridée}

Reprise du flux multi-D
$w\big[(U_p\otimes v_p - \tfrac{h_p}{2}\nabla_p U\cdot S_{\max})\cdot n\big] + (1-w)F_{1D}$
avec $S_{\max}=\mathrm{diag}(|v_{p,k}|)$ et $h_p = V_p/\sum_f A_f$, dans une routine
propre, la partie Lagrange (pression nodale) restant inchangée.

\begin{table}[h]\centering
\begin{tabular}{llcc}
\toprule
& hybridation & $L_2(\rho)$ à $\mathrm{Ma}=10^{-4}$ & flux quad \\
\midrule
\texttt{WIP2\_NOLM} & --- (advection 1D) & $\mathbf{6{,}03\cdot10^{-8}}$ & 6{,}98e-4 \\
\texttt{WIP2\_MD}   & $w=(\min A/A)/3$, façon AMISO & $1{,}26\cdot10^{-6}$ & 5{,}05e-4 \\
\texttt{WIP2\_MD5}  & $w=0{,}5$, façon ARMD & $1{,}12\cdot10^{-4}$ & --- \\
\texttt{WIP2\_MDH}  & idem MD, gradient enthalpique & $1{,}26\cdot10^{-6}$ & 4{,}92e-4 \\
\bottomrule
\end{tabular}
\end{table}

L'hybridation vaut \textbf{deux ordres de grandeur} ($1{,}26\cdot10^{-6}$ contre
$1{,}12\cdot10^{-4}$) : les \texttt{ARMD*} n'étaient pas mauvaises en soi, elles
étaient hybridées deux fois trop fort. Mais même bien dosée, l'advection multi-D
conserve un plancher bas Mach $20\times$ pire que l'advection 1D.

La correction enthalpique du gradient (\texttt{MDH}, $\nabla(\rho E)\to\bar h\nabla\rho$,
condition de Häenel appliquée au gradient nodal) est \textbf{numériquement sans effet}
($1{,}25636\cdot10^{-6}$ contre $1{,}25617\cdot10^{-6}$). Le plancher ne vient donc pas
de la ligne énergie mais des lignes densité et quantité de mouvement : dans ce
formalisme, tout terme dissipatif bâti sur $h_p\nabla U$ plutôt que sur un saut porte
ce risque, quelle que soit la variable dissipée.

\section*{5. Côté lagrangien}

Vortex de Gresho à $\mathrm{Ma}=10^{-2}$, ordre 1, énergie cinétique restante à $t=0{,}2$ :

\begin{table}[h]\centering
\begin{tabular}{lcc}
\toprule
& quad & tri \\
\midrule
\texttt{classic} (vitesse nodale) & $4\cdot10^{-5}$ & $0$ \\
\texttt{vp\_pp} (v et p par systèmes nodaux) & $\mathbf{1{,}012}$ & $1{,}151$ \\
\bottomrule
\end{tabular}
\end{table}

\texttt{classic} dissipe $99{,}996\,\%$ de l'énergie cinétique ; la dissipation y opère
à l'échelle \emph{acoustique} et non convective, donc $100\times$ trop vite à
$\mathrm{Ma}=10^{-2}$ (vérifié : à $\mathrm{Ma}=1$ il n'en perd que $31\,\%$).

\paragraph{Convergence en Mach, mesurée aussi côté lagrangien.} Même protocole qu'en
Euler ($t_{final}=10^{-2}$, $L_2(\rho)=\sqrt{\sum_i V_i(\rho_i-1)^2}$) :

\begin{table}[h]\centering
\begin{tabular}{lcccc}
\toprule
& \multicolumn{2}{c}{\texttt{classic}} & \multicolumn{2}{c}{\texttt{vp\_pp}} \\
$L_2(\rho)$ & quad & tri & quad & tri \\
\midrule
$\mathrm{Ma}=10^{-2}$ & 1{,}75e-6 & 3{,}23e-6 & 4{,}07e-7 & 1{,}53e-5 \\
$\mathrm{Ma}=10^{-4}$ & \textbf{2{,}31e-10} & \textbf{2{,}17e-10} & 3{,}89e-9 & 1{,}94e-6 \\
\midrule
$E_{kin}$ restante à $\mathrm{Ma}=10^{-4}$ & \textcolor{BrickRed}{$0{,}000$} & \textcolor{BrickRed}{$0{,}000$} & $1{,}0000$ & $1{,}0009$ \\
\bottomrule
\end{tabular}
\end{table}

\noindent
C'est la \textbf{seconde démonstration indépendante} du piège de la section~1, et la plus
propre : \texttt{classic} atteint $\mathrm{Ma}^2$ alors qu'il a intégralement détruit le
tourbillon \emph{avant} l'instant de mesure. Ici il n'y a aucun terme d'advection à
incriminer --- seul le couplage nodal agit.

\paragraph{L'énergie cinétique totale ne suffit pas non plus.} Elle est globale et
manque un défaut localisé. Bornes des champs à $t_{max}$, sur les mêmes runs :

\begin{table}[h]\centering
\begin{tabular}{lcc}
\toprule
 & $\rho$ & $|v|_{max}$ (physique : $1{,}0$) \\
\midrule
\texttt{classic} quad & $1{,}0000$--$1{,}0001$ & \textcolor{BrickRed}{$0{,}001$} \\
\texttt{classic} tri  & $1{,}0000$--$1{,}0000$ & \textcolor{BrickRed}{$0{,}000$} \\
\texttt{classic} Voronoï & $1{,}0000$--$1{,}0001$ & \textcolor{BrickRed}{$0{,}012$} \\
\midrule
\texttt{vp\_pp} quad & $0{,}9996$--$1{,}0004$ & $1{,}032$ \\
\texttt{vp\_pp} Voronoï & $0{,}9984$--$1{,}0018$ & $1{,}263$ \\
\texttt{vp\_pp} tri  & \textcolor{BrickRed}{$0{,}9045$--$1{,}1312$} & \textcolor{BrickRed}{$3{,}145$} \\
\bottomrule
\end{tabular}
\end{table}

\noindent
Deux lectures s'imposent. D'abord, la plage de densité \emph{impeccable} de
\texttt{classic} est la signature d'un \textbf{champ mort} et non d'une qualité :
$|v|_{max}$ y vaut $10^{-3}$ pour un maximum physique de $1$. Ensuite,
\texttt{vp\_pp} sur \texttt{tri} dépasse de \textbf{trois fois} ce maximum physique et
sa densité s'écarte de $13\,\%$, sans produire un seul NaN --- ce résultat passe donc
tous les contrôles habituels, et l'énergie cinétique totale ($1{,}0009$, soit
$+0{,}09\,\%$) ne le voit pas, le défaut étant confiné à quelques mailles.
\texttt{tri} est ainsi le cas dur pour le \emph{pas de temps} chez \texttt{classic}
(facteur $66$, section~5.1) et pour la \emph{qualité de la solution} chez
\texttt{vp\_pp}.

\medskip
\noindent
\texttt{vp\_pp} est la variante ajoutée ici : vitesse et pression nodales issues
\emph{toutes deux de systèmes nodaux}, donc sans aucune longueur $h_p$ ad hoc. Elle
s'interprète exactement comme \texttt{sidil} --- en développant, $p_p \approx \langle p\rangle
- \tfrac{h_p}{2}\rho_p a_p \nabla\!\cdot\!v$ et $u_p \approx \langle v\rangle -
\tfrac{h_p}{2\rho_p a_p}\nabla p$ --- à ceci près que $h_p$ y \emph{émerge} comme
$V_p/\sum_f A_f$ pondéré par les impédances réelles au lieu d'être choisi.
L'énergie totale est conservée à 6 chiffres ; la dérive de l'énergie cinétique
($+1{,}2\,\%$ quad, $+15\,\%$ tri) est un transfert depuis l'énergie interne de
$7\cdot10^{-7}$ en relatif, amplifié par le rapport $E_{int}/E_{kin}=2{,}1\cdot10^5$.

\subsection*{5.1 Les trois modes d'échec de \texttt{vp\_pp} au choc fort}

Sur le Sedov, \texttt{vp\_pp} échoue de \textbf{trois façons distinctes} selon le
maillage et l'ordre. Les confondre conduirait à chercher une cause unique là où il y
en a au moins deux.

\begin{table}[h]\centering
\begin{tabular}{llll}
\toprule
cas & symptôme & \texttt{EXIT\_CODE} & diagnostic fiable \\
\midrule
quad, tri (o1 et o2) & NaN par perte de positivité & \textcolor{BrickRed}{$0$} & comptage de NaN \\
Voronoï, ordre 2 & effondrement de $dt$, \emph{aucun} NaN & --- & suivi de $dt$ \\
Voronoï, ordre 1 & enchevêtrement de mailles, non fatal & $0$ & inspection du maillage \\
\bottomrule
\end{tabular}
\end{table}

\paragraph{NaN par perte de positivité (quad, tri).} Dans le quasi-vide, l'énergie
ambiante ($10^{-12}/(\gamma-1)$) est négligeable devant le travail reçu
($\sim\!10^{-8}$) : une maille fournit plus de travail qu'elle n'a d'énergie,
$\mathrm{sol}(5)$ devient négatif, et $\lambda=\max(\sqrt{\gamma p\tau}/\tau,\dots)$
prend la racine d'un argument négatif. Le terme acoustique est le seul du $\max$ à ne
pas être protégé. Le calcul diverge en moins de $100$ itérations et \textbf{sort avec
\texttt{EXIT\_CODE: 0}}, $t$ étant devenu NaN et \texttt{do while (t < t\_max)} étant
alors faux.

\paragraph{Effondrement du pas de temps (Voronoï, ordre 2).} Aucun NaN, aucune
grandeur non physique : $dt$ chute de cinq ordres entre $t=0{,}79$ et $t=0{,}898$
($2{,}98\cdot10^{-4} \to 3{,}10\cdot10^{-9}$) puis stagne, ce qui demanderait
$\sim\!3\cdot10^{7}$ itérations supplémentaires. Le champ à $t=0{,}80$ est sain. C'est
la signature d'une maille qui dégénère, rayon inscrit tendant vers zéro.

\paragraph{Ce que le fichier de sortie ne dit pas.} Il est écrit \emph{après} la
divergence et ne reflète pas l'état au moment de l'échec. Sur \texttt{tri} il ne
montre aucune pression négative alors qu'il en existe une, transitoire, dans une
maille qui finit NaN. Sur ces calculs, \textbf{ni les bornes du maillage ni les champs
finaux ne sont un diagnostic fiable}.

\subsection*{5.2 \texttt{vp\_pp} est meilleur que \texttt{classic} là où il tient}

Le défaut de positivité ne doit pas masquer l'inverse. Sur un maillage de Voronoï
\emph{radial} (disque $R=1{,}2$, anneaux concentriques, 5017 mailles, Sedov
$\varepsilon_0=0{,}984042$, CFL $0{,}1$, $t_{max}=1$, ordre 1, séquentiel) :

\begin{table}[h]\centering
\begin{tabular}{lccc}
\toprule
 & rayon du choc & dispersion & $\rho_{max}$ (exact $6{,}00$) \\
\midrule
\texttt{classic} & $0{,}982$ & $\pm1{,}3\,\%$ & $5{,}34$ \\
\texttt{vp\_pp}  & $0{,}991$ & $\mathbf{\pm0{,}2\,\%}$ & $6{,}65$ \\
\bottomrule
\end{tabular}
\end{table}

\noindent
\texttt{vp\_pp} est donc \textbf{six fois plus régulier sur le front}, et supprime
complètement les quatre encoches sur les axes que présente \texttt{classic}. Il
dépasse en contrepartie le pic exact.

\paragraph{Ce que cette dispersion ne mesure pas.} Elle vaut
$(\max-\min)/(2\,\mathrm{moy})$ du rayon du centroïde de densité maximale sur $36$
secteurs, c'est-à-dire du \textbf{bord externe} du choc. Or l'artefact en lobes vit
sur le \textbf{bord de la cavité}, et ces deux interfaces ne portent pas la même
information. Mesuré sur le nid d'abeilles, le mode azimutal du front a une amplitude
de $0{,}003$ --- indiscernable du bruit --- contre $0{,}10$ sur le bord de la cavité,
soit un facteur trente. \textbf{Une métrique bâtie sur le front ne quantifie donc pas
les lobes}, et peut annoncer une dispersion de quelques pour-cent sur un champ
visiblement lobé. Les chiffres ci-dessus restent valides pour ce qu'ils disent --- la
régularité du front, où l'écart entre les deux schémas est réel --- mais ils ne
mesurent pas l'amplitude de l'artefact.

\emph{Mise en garde sur ce maillage} : l'artefact à quatre lobes sur les axes est un
effet d'\textbf{ordre 1}, et non de la symétrisation employée pour construire le
disque --- un quart de disque seul présente le même défaut, et l'ordre 2 le supprime
dans les deux géométries.

\medskip
\noindent
Ces trois modes d'échec et ces mesures radiales proviennent du travail de la session
en charge du module lagrangien. Aucun correctif n'a été appliqué : le limiteur de
positivité et la protection du terme acoustique sont des modifications de schéma, en
attente d'arbitrage.

\section*{Conclusions}

\begin{enumerate}\itemsep3pt
  \item \textbf{Le critère~2 ne se lit jamais seul} : quatre schémas le « gagnent » en
        ayant détruit le vortex. Toujours l'accompagner du critère~1.
  \item \textbf{\texttt{WIP2\_NOLM} est le seul des sept schémas de référence à passer
        les quatre critères}, et son ingrédient décisif est la vitesse d'énergie
        symétrique (conservative) associée à la pression nodale : à pression nodale
        identique, la forme unilatérale de \texttt{LPP} plafonne à
        $3{,}2\cdot10^{-4}$ contre $6{,}0\cdot10^{-8}$.
  \item \textbf{Le terme $\epsilon$ ne peut pas être supprimé}, ni en le reformulant
        sans capteur, ni en changeant d'advection. Le capteur de Ducros en place est
        le meilleur mesuré.
  \item \textbf{La tension quad/tri est un effet d'amplitude}, monotone et de signe
        opposé sur les deux maillages. C'est, à mon sens, la question centrale
        restante : elle est apparue quatre fois dans cette campagne par des mécanismes
        indépendants.
\end{enumerate}

\paragraph{Réserves.} L'échelle absolue du nombre de Stanton provient de la formule déjà
présente dans \texttt{plot\_st.gnu} et n'a pas été re-dérivée : les classements sont
fiables, la valeur absolue mériterait vérification. Les maillages triangulaires
dépendent de la version de \texttt{gmsh} ; seules les comparaisons à maillage identique
ont un sens, ce qui est le cas ici. Enfin \texttt{ZB\_ARMDWIP\_*} reste sélectionnable
alors que \texttt{mat\_h\_p} est nul depuis le retrait de \texttt{compute\_ellip} : il
produirait des résultats silencieusement faux et devrait être désenregistré ou réparé.

\end{document}
